\documentclass[12pt]{article}
\usepackage[margin=0.75in]{geometry}

\newcommand{\duedate}{01/28/2026}
\newcommand{\assignment}{Continual Learning Overview: Summary}

\newcommand{\name}{Aksel Kretsinger-Walters, adk2164}
\newcommand{\email}{adk2164@columbia.edu}

\makeatletter
\def\input@path{{../}{../../}{../../../}}
\makeatother
\input{pset_template_CLMM.tex} %% DO NOT CHANGE THIS LINE

\begin{document}
\psetheader %% DO NOT CHANGE THIS LINE

\section*{Nested Learning}

\paragraph{Motivation.}
Modern machine learning systems combine neural architectures and optimization algorithms
in complex ways. Despite advances in deep learning, there remains limited understanding
of how learning algorithms, architectures, and memory mechanisms interact.

The Nested Learning (NL) paradigm proposes a unifying framework that models learning
systems as hierarchies of interconnected optimization processes.

\paragraph{Core Idea.}
Nested Learning represents a machine learning model as a collection of
\textbf{nested optimization problems}, each operating at a different level
of abstraction. \cite{nestedlearning}

In this view:

\begin{itemize}
		\item Neural architectures operate on \textbf{data context} (tokens, images, etc.).
		\item Optimizers operate on \textbf{gradient context}.
\end{itemize}

Both processes can be interpreted as learning systems that compress their respective contexts.

This perspective suggests that optimization algorithms and neural architectures
are fundamentally the same type of object operating at different levels.

\paragraph{Associative Memory Interpretation.}
Nested Learning interprets many machine learning components as forms of associative memory.

For example:

\begin{itemize}
		\item Neural networks compress training data into parameters.
		\item Optimizers compress gradient information into momentum or adaptive statistics.
\end{itemize}

Under this interpretation, common optimizers such as Adam or SGD with momentum
are themselves associative memory systems that learn to store gradient information.

\paragraph{Multi-Level Learning Systems.}
The NL framework introduces multiple interacting learning levels:

\begin{enumerate}
		\item Data-level learning (standard neural network training)
		\item Optimizer-level learning (learning how to update parameters)
		\item Meta-learning levels that adapt the learning algorithm itself
\end{enumerate}

Each level maintains its own internal gradient flow and update rule.

\paragraph{Self-Modifying Models.}
Building on this framework, the paper proposes architectures that can modify
their own learning rules. These models learn update rules that control how
their internal memory and parameters evolve over time.

This enables forms of continual learning and adaptive reasoning.

\paragraph{Continuum Memory System.}
The paper also introduces a \emph{continuum memory system} in which memory
exists along a spectrum of update frequencies rather than being divided into
discrete short-term and long-term components.

Different memory modules update at different temporal scales, enabling
both rapid adaptation and persistent knowledge storage.

\paragraph{Hope Architecture.}
As a proof of concept, the authors introduce a continual learning model
called \emph{Hope}, which integrates:

\begin{itemize}
		\item nested optimization mechanisms
		\item adaptive memory modules
		\item self-modifying learning rules
\end{itemize}

Experiments demonstrate promising results on language modeling,
continual learning benchmarks, and few-shot generalization tasks.

\paragraph{Limitations.}
The Nested Learning framework is primarily conceptual and theoretical.
Many proposed mechanisms remain difficult to implement efficiently
in large-scale models.

Additionally, the framework introduces substantial architectural
complexity and raises challenges for stable training.

\paragraph{Takeaway.}
Nested Learning proposes a new perspective on deep learning in which
architectures, optimizers, and memory systems are viewed as nested
learning processes. This unified framework may provide insights for
designing more adaptive and continual learning systems.

\end{document}
